\section{Details on crowdworker recruitment, data labeling services, quality control} \label{app:quality_control}

\paragraph{Fair and Ethical Labor} We have enlisted the full-time services of 70 crowdworkers, notable for their proficiency in text annotation for commercial machine learning projects, and their ability to navigate multifaceted issues such as determining risk neutrality between pairs of harmful prompts and harmless responses. In recognition of their valuable contributions, we have established an equitable compensation structure. Their estimated average hourly wage ranges from USD 7.02 to USD 9.09 (XE rate as of 2023/06/07), significantly exceeding the minimum hourly wage of USD 3.55~\cite{noauthor_undated-nz} (XE rate as of 2023/06/07) in Beijing, PRC. In adherence to local labor laws and regulations, our crowdworkers follow an established work schedule consisting of eight-hour workdays on weekdays, with rest periods during weekends. 

\vspace{-.7em}
\paragraph{Fair Use of Dataset and Identifying Potential Negative Societal Impacts} The \textsc{BeaverTails} project has undergone thorough review and auditing by the Academic Committee of the Institution for Artificial Intelligence at Peking University. The committee has served as the Institutional Review Board (IRB) for this work and ensures that the use of the \textsc{BeaverTails} dataset adheres to principles of fairness and integrity.

\paragraph{Lessons Learned: Identifying harmlessness for a QA pair is a complex, multi-faceted task}

During the initial fortnight of our project, we utilized a single-stage annotation model where crowdworkers first assessed the safety of the QA pair and then ranked responses by their helpfulness and harmlessness in one attempt. However, this model presented considerable alignment difficulties between the research and annotation teams, particularly around defining what constitutes a harmless response when faced with a harmful prompt. Significant disagreements arose around the harmlessness of a QA pair, leading to a large portion of preference-ranking data being rendered unusable. This issue was largely due to the premise that ranking data only holds meaningful value when the safety of a QA pair is accurately labeled. These disagreements were particularly pronounced in two key areas: the degree to which a response's correctness should be tied to its harmlessness and how an AI assistant should approach sensitive topics such as marijuana legalization or gun control. We realized that these issues resulted from overly simplistic criteria for categorizing a QA pair as harmful or harmless. To address this, we reconsidered our approach and adopted a two-stage model that separated the complex, multi-faceted task of identifying harmfulness into discerning the presence of 14 specific potential harm categories. This shift led to an approximately 15\% increase in agreement rates during our quality control tests, indicating an improved alignment between the researchers and annotators.

\paragraph{Recruiting crowdworkers and data-labeling service provider}
We have collaborated with a professional data annotation service provider called~\href{www.aijetdata.com}{AIJet Data}. We did not directly engage with the crowdworkers; AIJet took charge of this process. Given AIJet's expertise in text-based data annotation, they assembled a team of skilled data annotators for our project. Recognizing the project's complexity, we agreed to a contract priced above the standard market rate, enabling us to prioritize the qualifications of the annotators. All chosen annotators were proven to have successfully completed the College English Test. Beyond this, they underwent a rigorous screening process, requiring them to achieve at least 90\% accuracy on a test aligned with our research team's answers. As a result, our team of 70 members was selected after a pool consisting of roughly 200 people. Only after passing this test were they formally contracted. We have provided them with a comprehensive annotation guideline to ensure adherence to our standards (Appendix ~\ref{app:annoation-documents}).

\paragraph{Quality Control Process} The quality control (QC) process we follow operates roughly in this manner:
\begin{itemize}[leftmargin=*]
% \setlength\itemsep{-1em}
    \item Three entities participate in the QC process: the data annotators, the AIJet QC team, and our research team. The AIJet team manages the assignment of workloads, the training of workers, and the collection of questions from the workers, which are then discussed with the research team (which occurred almost daily between April and May).
    \item Once a data annotator completes an assigned batch, the internal system forwards this batch to the AIJet QC team. The AIJet QC team members review each annotated pair based on the standards set by the research team. The inspected batch is then forwarded to the research team for additional quality assessments. According to our agreed terms, we must sample at least 10\% of the data from the inspected batches, and the percentage agreement must reach a minimum of 90\% for acceptance. We set this threshold because achieving 100\% agreement is not realistically feasible, nor is it commercially viable for the data annotation service provider. It also runs the risk of introducing further biases from the research team. For a batch to be rejected, at least two research team members must inspect it.
    \item The initial stages of our collaboration with AIJet were quite challenging. During the first two weeks, we rejected all of their inspected batches, prompting AIJet to urgently request several face-to-face meetings with the research team. Over two months, the agreement rate gradually climbed from the 60\%-70\% range to the 88\%-92\% range. A significant factor contributing to this improvement was the introduction of the two-stage annotation model. We found that breaking down our rigid standards into a series of binary choice questions greatly assisted the data annotators in understanding our intentions.
\end{itemize}



